\lead{Hyperball needs a large LR and still overfits.} Hyperball projects every Linear weight back to its initial Frobenius norm after each Adam step and uses WD 0, so the norm cannot grow and the only free knob is how fast the direction turns. On this 32-pass DNA task it behaves very differently from AdamW with WD 3.2:
\begin{itemize}
\item At LR .001 and .003 the validation NLL falls to about 1.279 by pass 6 and then climbs steadily to 1.44 and 1.58 by pass 32 (the LR .003 run, which fits faster, ends worse), the same late-epoch failure as AdamW at WD .1 (1.4057).
\item The higher LR .01 delays the turn (best 1.268 at pass 14) but ends at 1.363. Larger LRs help because a norm-constrained step rotates the weights more, which acts like a stronger regularizer early on, but no LR in our range keeps the model on the floor for 32 passes. Fixing the weight norm removes weight growth, not the ability to fit 30k sequences many times over, so Hyperball at WD 0 is not a substitute for weight decay when the data are repeated this often.
\item The best early-stopped Hyperball value (1.268 to 1.279 at 6 to 14 passes) is within 0.01 of the AdamW baseline, which says the architecture and data, not the optimizer, set the floor; the optimizer's job here is to stop the model from leaving it.
\end{itemize}
